\subsection{时间余圈、交叉积寄存器与 Watatani 指数的可消去判据}\label{subsec:conclusion-time-cocycle-register-eliminability-index}
沿用正文一参数时间--模量框架中的重整化缺陷
\[
\Delta_m(w):=\ell_{m+1}(w)-\ell_m\!\bigl(p_{m+1,m}(w)\bigr),
\]
由长度余圈诱导的模流与交叉积时间寄存器，以及条件期望 \(E\) 的点态 Watatani 指数 \(\Ind(E)(x)\) 的记号。前文已经分别证明：跨分辨率缺陷的 \(w\)-依赖恰是无法由层间重标定吸收的一阶障碍；时间 gauge 自由度与 Connes \(1\)-余圈完全同一；时间寄存器正是交叉积变量；而局部最小辅助时间纤维维数则由 \(\Ind(E)(x)\) 精确刻画。把这四条链并读之后，可把“时间寄存器是否真能被消去”压缩为一个上同调--指数联合判据。

\begin{theorem}[时间寄存器可消去的上同调--指数判据]\label{thm:conclusion-time-register-eliminability-cohomology-index-criterion}
在上述一参数时间--模量框架中，设 \(E:\mathcal A\to\mathcal B\) 为与可逆提升实现对应的忠实正规条件期望。则下列两条陈述等价：
\begin{enumerate}
  \item 对每个分辨率层 \(m\)，重整化缺陷 \(\Delta_m(w)\) 对 \(w\) 为常数，且
  \[
  \Ind(E)(x)=1
  \qquad\text{a.e.};
  \]
  \item 时间 \(1\)-余圈在规范意义下同调平凡，因而交叉积时间寄存器可被完全消去；等价地，系统可在不引入任何非平凡辅助时间纤维的前提下实现因果可逆化。
\end{enumerate}

相反，若存在某层 \(m\) 使 \(\Delta_m(w)\) 具有非平凡 \(w\)-依赖，或存在正测度集合上
\[
\Ind(E)(x)>1,
\]
则每一因果可逆实现都必须携带非平凡时间寄存器，且其点态最小附加时间开销满足
\[
\boxed{\;
T_{\min}(x)\ \ge\ \frac{\log \Ind(E)(x)}{\log |A_0|}\; }.
\]
\leanverified{paper\_conclusion\_time\_register\_eliminability\_cohomology\_index\_criterion}
\end{theorem}
\begin{proof}
先证 \((1)\Rightarrow(2)\)。若每个 \(\Delta_m(w)\) 都仅依赖于 \(m\)，则定理 \ref{thm:xi-cross-resolution-time-renormalization-defect} 表明这些缺陷可被吸收进层间常数重标定，从而跨分辨率时间实现严格相容。再由定理 \ref{thm:xi-modular-flow-equals-time-cocycle}，这正意味着时间 \(1\)-余圈与一个 gauge 余边界同调等价，即其 Connes 类平凡。最后，定理 \ref{thm:xi-time-register-crossed-product-equivalence} 识别出时间寄存器恰是交叉积变量，因此该同调平凡性正等价于时间寄存器可被消去。条件 \(\Ind(E)(x)=1\) a.e. 则说明局部时间纤维上没有任何剩余 multiplicity 障碍。

再证 \((2)\Rightarrow(1)\)。若时间寄存器可被消去，则不存在非平凡中心扩张需要外置，故由定理 \ref{thm:xi-cross-resolution-time-renormalization-defect}，每个 \(\Delta_m(w)\) 必对 \(w\) 为常数。另一方面，寄存器可消去意味着局部最小辅助时间纤维维数处处为 \(1\)；再由定理 \ref{thm:xi-local-index-equals-local-time-fiber}，
\[
\Ind(E)(x)=1
\qquad\text{a.e.}
\]
随之成立。

最后看阻碍项。若某层 \(\Delta_m\) 具有非平凡 \(w\)-依赖，则定理 \ref{thm:xi-cross-resolution-time-renormalization-defect} 已断言：任何全局时间一致化都需要额外中心扩张，故时间寄存器不可消去。若 \(\Ind(E)(x)>1\) 在正测度集合上成立，则由定理 \ref{thm:xi-local-index-equals-local-time-fiber}，该集合上的局部最小辅助时间纤维维数严格大于 \(1\)，从而任何因果可逆实现都必须携带非平凡时间寄存器，并满足
\[
T_{\min}(x)\ \ge\ \frac{\log \Ind(E)(x)}{\log |A_0|}.
\]
证毕。
\end{proof}

\leanverified{paper\_conclusion\_anomaly\_strictification\_universal\_time\_host}
\begin{theorem}[异常吸收到加性时间寄存器的普适 strictification]\label{thm:conclusion-anomaly-strictification-universal-time-host}
在上述一参数时间--模量框架中，设某个 strictification 把异常/反项数据吸收到一条加性时间寄存器中，并保持协变表示与时间坐标的加法结构。
则该 strictification 必唯一因子化经由半直积群胚与交叉积对象
\[
\boxed{\;
W^\ast(\widetilde{\mathcal G})
\cong
\mathcal M\rtimes_{\sigma^{\varphi_\lambda}}\mathbb R,\qquad
\widetilde{\mathcal G}:=\mathcal G\ltimes_{\widetilde\ell}\mathbb R\; }.
\]
并且对每个中心谱点 \(x\)，所需最小辅助时间纤维维数满足
\[
\boxed{\;
\dim(\mathcal T_x)_{\min}=\Ind(E)(x).\;}
\]
因此，在保持协变性与时间加法结构的前提下，时间轴不是任意附加的解释变量，而是异常 strictification 的普适且点态极小宿主。
\end{theorem}
\begin{proof}
由定理 \ref{thm:pom-strictification-splitting}，异常消去等价于为相应中心扩张选择一条分裂；换言之，任何 strictification 都是在某个加法寄存器上实现该分裂。
另一方面，定理 \ref{thm:xi-time-register-crossed-product-equivalence} 已证明：凡把异常吸收到时间中心扩张并保持时间加法结构的实现，统一落在同一个普适对象
\(
W^\ast(\widetilde{\mathcal G})\cong \mathcal M\rtimes_{\sigma^{\varphi_\lambda}}\mathbb R
\)
上，故上述 strictification 必唯一因子化经由该交叉积。
最后，定理 \ref{thm:xi-local-index-equals-local-time-fiber} 直接给出局部最小辅助时间纤维维数
\(
\dim(\mathcal T_x)_{\min}=\Ind(E)(x)
\)。
三者合并即得结论。证毕。
\end{proof}

\paragraph{统一读法}
这一判据说明：时间寄存器并不是实现细节，而是由两个彼此独立的障碍共同决定的最小外置变量。一端是跨分辨率缺陷的上同调平凡性，另一端是局部时间纤维的指数型 multiplicity。只有当二者同时塌缩到平凡值时，时间寄存器才真正可消去。

\endinput
